\chapter{Conclusion}\label{ch:conclusion}

\section{Result Discussion and Conclusion}
\subsection{Peak Detection Algorithm}
We developed a basic algorithm model for identifying QRS complexes based upon a Processing Phase, and then 
Identification Phase using a adaptive threshold, we tested the performance of this algorithm against the QRS annotation 
files, noting that different time interval thresholds worked for different datasets. We then produced extensions of this algorithm to improve 
the performance, concluding with a Jump Detection Based Algorithm that produces the best performance, at the sacrifice 
of time complexity. 
\subsection{Respiratory Rate Methods}
We evaluated the RSA and RPA Respiratory Rate estimation methods, using spectrograms as the predomiant tool with the goal 
of identifying a dominate frequency, that persitted across time. We began by evaluating two methods 
for spectrum analysis of the non-uniform time series, a Discrete Fourier Transform with interpolation and LombScargle. 
We saw that the interpolation was very accurate, and due to the error involved in the LombScargle method along 
with time complexity, that the FFT based method was better for producing Spectrograms. We saw that the RPA method 
produced a stronger dominate frequency when it was able to detect one, however the RSA method was more robust across 
the datasets, however neither performed very well at all, with 7 of the 8 tested datasets not having a visually obvious pattern 
in the Spectrogram. 
\section{Area for future research}
We suggeset the following directions for extending the work of this Project. \\
Peak Detection Algorithm:
\begin{itemize}
        \item Implement statistical methods to determine false detection or missed detection of 
        RR peaks, by identifying small or large gaps in RR peak intervals based off the average interval size. 
        When detected remove, or adjust threshold. 
\end{itemize}
Respiratory Rate Methods:
\begin{itemize}
        \item Adjust Spectrogram increments and window sizes to see if this improves ability 
        to detect signals.
        \item Estimate Error of the interpolation method to provide further validity of this FFT based 
        method.
        \item Testing Methods with other known QRS complex algorithms, e.g Pan Tompkins.
        \item Implement unsupervised learning techniques to identify patterns in the Spectrogram that may 
        correspond to Respiratory Rate that aren't visually apparent. 
\end{itemize}